\documentclass[11pt]{article}
\usepackage[a4paper,margin=1in]{geometry}
\usepackage[T1]{fontenc}
\usepackage{lmodern}
\usepackage{microtype}
\usepackage{xcolor}

\setlength{\parindent}{0pt}
\setlength{\parskip}{0.5em}

\newcommand{\promptbox}[1]{%
  \begingroup\setlength{\fboxsep}{10pt}\noindent
  \colorbox{blue!8}{\parbox{\dimexpr\linewidth-2\fboxsep\relax}{\textbf{Prompt. }#1}}
  \endgroup
}
\newcommand{\resultbox}[1]{%
  \begingroup\setlength{\fboxsep}{10pt}\noindent
  \colorbox{purple!8}{\parbox{\dimexpr\linewidth-2\fboxsep\relax}{\textbf{Result. }#1}}
  \endgroup
}
\newcommand{\examplebox}[1]{%
  \begingroup\setlength{\fboxsep}{10pt}\noindent
  \colorbox{green!8}{\parbox{\dimexpr\linewidth-2\fboxsep\relax}{\textbf{Example. }#1}}
  \endgroup
}

\title{\textbf{Book 1 --- Lecture 1 --- Exercise 1}\\[2mm]
\large Closed vs Open Systems}
\author{}
\date{}

\begin{document}
\maketitle

\promptbox{Since the notion is so important to theoretical physics, think about
what a closed system is and speculate on whether closed systems can actually
exist. What assumptions are implicit in establishing a closed system? What is an
open system?}

\section*{Closed System (working definition)}
In theoretical physics, a \emph{closed} system is one you describe with a
\emph{complete} set of dynamical variables and \emph{deterministic} (or, in
quantum theory, unitary) laws that involve \emph{only} those variables.
Nothing outside the model is allowed to influence the state except through
quantities you have already included.\footnote{A \emph{dynamical variable} is a
quantity whose value changes with time and helps specify the system state. A law
is \emph{deterministic} when the current state fixes a unique next state.}

\section*{Implicit Assumptions}
Calling a system ``closed'' usually includes several idealizations:
\begin{itemize}
  \item \textbf{Exhaustive state:} all relevant degrees of freedom are modeled.\footnote{A
  \emph{degree of freedom} is an independent parameter needed to describe a
  system configuration.}
  \item \textbf{Known laws:} dynamics is assumed exact for that state.
  \item \textbf{Isolation:} outside influences are negligible or explicitly modeled.
  \item \textbf{Boundary in time and space:} closure is only approximate over a chosen window.
\end{itemize}

\section*{Can Closed Systems Actually Exist?}
Operationally, we can approximate closure very well. Fundamentally, perfect
closure is doubtful: long-range gravity, environmental coupling, and quantum
entanglement are hard to eliminate completely. In practice, ``closed'' is a
controlled approximation whose quality depends on the required precision.

\section*{Open System}
An \emph{open} system exchanges energy, momentum, particles, or information with
an environment that is not fully tracked. Its effective dynamics then includes
dissipation, noise, driving, or explicit coupling terms. Often the combined
``system + environment'' is treated as closed, while the subsystem of interest
is open.\footnote{A \emph{subsystem} is a selected part of a larger physical
system. \emph{Dissipation} means irreversible energy transfer (often to heat).}

\resultbox{Classical mechanics often starts with closed deterministic models.
When that is insufficient, one either enlarges the modeled system or uses
open-system descriptions.}

\examplebox{A pendulum in vacuum over short times can be modeled as nearly
closed; adding air drag and thermal fluctuations turns it into an open-system
description.}

\end{document}
